\documentclass[12pt,a4paper]{article}

% 基本宏包
\usepackage[utf8]{inputenc}   % UTF-8 编码
\usepackage[T1]{fontenc}      % 更好的英文字体输出
\usepackage{lmodern}          % Latin Modern 字体
\usepackage{amsmath, amssymb, amsthm} % 数学环境
\usepackage{geometry}         % 控制页边距
\usepackage{hyperref}         % 超链接
\usepackage{enumitem}         % 列表控制
\usepackage{graphicx}         % 插图
\usepackage{tcolorbox}        % 彩色方框，做重点标注
\usepackage{fancyhdr}         % 页眉页脚
\usepackage{braket}           % Dirac notation: \ket, \bra, etc.

\numberwithin{equation}{subsection}

% 页边距
\geometry{margin=1in}

% 页眉页脚
\pagestyle{fancy}
\fancyhf{}
\lhead{Quantum field theory}   % 左页眉
\rhead{\today}     % 右页眉
\cfoot{\thepage}   % 页脚居中页码

% 彩色框样式
\tcbset{colback=blue!5!white, colframe=blue!75!black, boxrule=0.8pt, arc=4pt, left=6pt, right=6pt, top=6pt, bottom=6pt}

% 定理环境
\newtheorem{theorem}{Theorem}
\newtheorem{definition}{Definition}

\begin{document}

\title{Learning Notes (Quantum field theory)}
\author{Bing}
\date{\today}
\maketitle

\tableofcontents  % 目录
\newpage

\section{Spin Zero}

\subsection{Attempts at relativistic quantum mechanics}
The axiom we need to focus on is the one that says the time
evolution of the state of the system is governed by the
Schrödinger equation
\begin{equation}
i\hbar\frac{\partial}{\partial t}|\psi,t\rangle =\hat{H}|\psi,t\rangle
\end{equation}
A spinless,nonrelativistic particle with no forces acting on it.
The hamiltonian is
\begin{equation}
H = \frac{1}{2m}\mathbf{P}^2
\end{equation}
Eq.(1.1.1) becomes
\begin{equation}
i\hbar\frac{\partial}{\partial t}\psi(\mathbf{x},t)=-\frac{\hbar^2}{2m}\nabla^2\psi(\mathbf{x},t)
\end{equation}
Generalize this to relativistic motion.
\begin{equation}
H=+\sqrt{\mathbf{P}^2c^2+m^2c^4}
\end{equation}
Expand this hamiltonian in inverse powers of the speed of light c.
\begin{equation}
H=mc^2+\frac{1}{2m}\mathbf{P^2}+\dots
\end{equation}
The Schrödinger becomes
\begin{equation}
i\hbar\frac{\partial}{\partial t}\psi(\mathbf{x},t)=+\sqrt{-\hbar^2c^2\nabla^2+m^2c^4}\psi(\mathbf{x},t)
\end{equation}
Difficulties:
\begin{itemize}
    \item It apparently treats treats space and time on difference footing.
    \item If we expand the square root in powers of $\nabla^2$, we get an infinite number of spatial
derivatives acting on $\psi(\mathbf{x},t)$ which implies that eq.(1.1.6) is not local in space.
\end{itemize}
Squaring the differential operators on each side of eq.(1.1.6).\\

\textit{Klein-Gordon equation}
\begin{equation}
-\hbar^2\frac{\partial^2}{\partial t^2}\psi(\mathbf{x},t)=(-\hbar^2c^2\nabla^2+m^2c^4)\psi(\mathbf{x},t)
\end{equation}
It is secomd-order in both space and time derivatives,
and they appear in a symmetric fashion.\\
Special relativity tells us taht physical looks like the same in all interial frames.
Suppose that a certain spacetime coordinate system $(ct,\mathbf{x})$ represents an inertial frame.
Define:
$$
x^0=ct,x^\mu\text{ in place of} (ct,\mathbf{x})\quad\mu=0,1,2,3
$$
$$
x_0=-x^0,x_i=x^i
$$
The Minkowsiki netric:
\begin{equation}
g_{\mu\nu}=
\begin{pmatrix}
-1&&&\\
&+1&\\
&&+1&\\
&&&+1
\end{pmatrix}
\end{equation}

$$
x_\mu=g_{\mu\nu}x^\nu
$$

\begin{equation}
g_{\mu\nu}=
\begin{pmatrix}
+1&&&\\
&-1&\\
&&-1&\\
&&&-1
\end{pmatrix}
\end{equation}
$$
x^\mu=g^{\mu\nu}x_\nu
$$

$$
g^{\mu\nu}g_{\nu\rho}=\delta^\mu_\rho\text{ (Kronecker delta)}
$$
Coordinates $x^\mu$ represent an intertial frame,then so do any other
coordinates $\overline{x}^\mu$ that are related by 
\begin{equation}
\overline{x}^\mu=\Lambda^\mu_\nu x^\nu+a^\mu
\end{equation}
where,
\begin{itemize}
    \item $\Lambda^\mu_\nu$ is a \textit{Lorentz transformation matrix}.
    \item $a^\mu$ is a translation vector.
    \item Both of two are constant.
\end{itemize}
$\Lambda^\mu_\nu$ must obey 
\begin{equation}
g_{\mu\nu}\Lambda^\mu_\rho\Lambda^\nu_\sigma =g_{\rho\sigma}
\end{equation}
The interval between two different spacetime points that are labeled by
$x^\mu$ and $x'^\mu$ in one other frame.
The interval:
\begin{equation}
    \begin{aligned}
(x-x')^2&=g_{\mu\nu}(x-x')^\mu(x-x')^\nu\\
&=(\mathbf{x}-\mathbf{x}')-c^2(t-t')^2
    \end{aligned}
\end{equation}
In other frame,
\begin{equation}
    \begin{aligned}
    (\overline{x}-\overline{x}')^2 &= g_{\mu\nu}(\overline{x}-\overline{x}')^\mu(\overline{x}-\overline{x}')^\nu\\
    &=g_{\mu\nu}\Lambda^\mu_\rho\Lambda^\nu_\sigma()(x-x')^\rho(x-x')^\sigma\\
    &=g_{\rho\sigma}(x-x')^\rho(x-x')\sigma\\
    &=(x-x')^2
    \end{aligned}
\end{equation}
Useful notation for spacetime derivates:
\begin{align}
\partial_\mu &=\frac{\partial}{\partial x^\mu}=(+\frac{1}{c}\frac{\partial}{\partial t},\nabla)\\
\partial^\mu &=\frac{\partial}{\partial x_\mu}=(-\frac{1}{c}\frac{\partial}{\partial t},\nabla)
\end{align}
\begin{equation}
\partial^\mu x^\nu=g^{\mu\nu}
\end{equation}
Because of eq.(1.10),
\begin{equation}
\overline{\partial}^\mu=\Lambda^\mu_\nu\partial^\nu
\end{equation}
\begin{equation}
\overline{\partial}^\rho\overline{x}^\sigma=(\Lambda^\mu_\nu\partial^\nu)(\Lambda^\sigma_\nu x^\nu+a^\mu)=\Lambda^\rho_\mu\Lambda^\sigma_\nu(\partial^\mu x^\nu)=\Lambda^\rho_\mu\Lambda^\sigma_\nu g^{\mu\nu}=g^{\rho\sigma}
\end{equation}
Write eq.(1.1.7) as 
\begin{equation}
-\hbar^2c^2\partial^2_0\psi(x)=(-\hbar^2c^2\nabla^2+m^2c^4)\psi(x)
\end{equation}
Rearrange and identify $\partial^2=\partial^\mu\partial_\mu=-\partial^2_0+\nabla^2$,
\begin{align}
(-\partial^2+\frac{m^2c^2}{\hbar^2})\psi(x)&=0\\
(-\partial^2+\frac{m^2c^2}{\hbar^2})\overline{\psi}(\overline{x}) &=0
\end{align}
We set $\overline{\psi}(\overline{x})=\psi(x)$,
\begin{equation}
\overline{\partial}^2=g_{\mu\nu}\overline{\partial}^\mu\overline{\partial}^\nu=g_{\mu\nu}\Lambda^\mu_\rho\Lambda^\nu_\sigma\partial^\rho\partial^\sigma=\partial^2
\end{equation}
Detailed proof:
\begin{enumerate}
    \item In the new reference frame, the d'Alembert operator is defined as:
    \[\partial^2=\partial^\mu\partial_\mu=g_{\mu\nu}\partial^\mu\partial^\nu
    \]
    \[
    \overline{\partial}^2 = g_{\mu\nu} \overline{\partial}^\mu \overline{\partial}^\nu
    \]

    \item The derivative operator transforms under Lorentz transformations by the inverse vector transformation:
    \[
    \overline{\partial}^\mu = \Lambda^\mu_\rho \partial^\rho, \quad 
    \overline{\partial}^\nu = \Lambda^\nu_\sigma \partial^\sigma
    \]
    \[
    \overline{\partial}^2 = g_{\mu\nu} (\Lambda^\mu_\rho \partial^\rho) (\Lambda^\nu_\sigma \partial^\sigma) 
    = g_{\mu\nu} \Lambda^\mu_\rho \Lambda^\nu_\sigma \partial^\rho \partial^\sigma
    \]

    \item Utilizing the metric invariance of the Lorentz transformation:
    \[
    g_{\mu\nu} \Lambda^\mu_\rho \Lambda^\nu_\sigma = g_{\rho\sigma}
    \]
    \[
    \overline{\partial}^2 = g_{\rho\sigma} \partial^\rho \partial^\sigma
    \]

    \item Finally, we notice
    \[
    g_{\rho\sigma} \partial^\rho \partial^\sigma = \partial_\sigma \partial^\sigma = \partial^2
    \]
    Therefore：
    \[
    \overline{\partial}^2 = \partial^2
    \]
\end{enumerate}
Therefore,The Klevin-Gordon equation is manifestly consistent with relativeity:it takes the same form in every inertial frame.
But,the abstract Schrödinger equation has the fundamental property of being first order in the time derivatve, whereas the Klevin-Gordon
equation is second order.

The norm of the state is not in general time independent.
\begin{equation}
\langle\psi,t|\psi,t\rangle=\int d^3x\langle\psi,t|\mathbf{x}\rangle\langle\mathbf{x}|\psi,t\rangle=\int d^3x\psi^*(x)\psi(x)
\end{equation} 
Introduce an extra discrete label on the wave function to account for spin:$\psi_a(x),a=1,2$.

Diac equation:
\begin{equation}
i\hbar\frac{\partial}{\partial t}\psi_a(x)=\left(-i\hbar c(\alpha^j)_{ab}\partial_j+mc^2(\beta)_{ab}\right)\psi_b(x)
\end{equation}
where all repeated indices are summed,and $\alpha^j$ and $\beta$ are matrices in spinspace.

The state $|\psi,a,t\rangle$ carries a spin label $a$, and the hamiltonian is
\begin{equation}
H_{ab}=c\hat{P_j}(\alpha^j)_{ab}+mc^2(\beta)_{ab}
\end{equation}
where $P_j$ is a component of the momentrum operator.

Square the hamiltonian
\begin{equation}
(H^2)_{ab}=c^2P_jP_k(\alpha^j\alpha^k)_{ab}+mc^3P_j(\alpha^j\beta+\beta\alpha^j)_{ab}+(mc^2)^2(\beta^2)_{ab}
\end{equation}
Since $P_jP_K$ is symmetric on exchange of $j$ and $k$, we can replace $\alpha^j\alpha^k$
by its symmetric part, $\frac{1}{2}\{\alpha^j,\alpha^k\}$,where $\{A,B\}=AB+BA$ is the anticommutator.
\begin{equation}
\{\alpha^j,\alpha^k\}_{ab}=2\delta^{jk}\delta_{ab},\quad\{\alpha^j,\beta\}_{ab}=0,\quad (\beta^2)_{ab}=\delta_{ab}
\end{equation}
\begin{equation}
(H^2)_{ab}=(\mathbf{P}^2c^2+m^2c^4)\delta_{ab}
\end{equation}
Thus, the eigenstates of $H^2$ are momentum eigenstates, with $H^2$ eigenvalue $\mathbf{p}^2c^2 + m^2c^4$.

We would like the Dirac matrices to be $2\times2$, in order to account for alectron spin.But they must in fact be larger and we find that the trace of this matrix is zero.Thus the
four eigenvalues must be $+E(\mathbf{p}),+E(\mathbf{p}),-E(\mathbf{p}),-E(\mathbf{p})$.The negative eigenvalues indicate that is no gound state which is a serious problem.
Diac sea of electrons: all the negative energy states were already occupied.In this case, a positive
energy electron could not drop into one of these states.The concept of a hole was interpreted as a positron, predicting the existence of antimatter.

\textbf{Limitations:}
\begin{itemize}
    \item A single relativistc particle to a theory that apparently requires an infinite number of particles.
    \item How can we solve the problem of how to describe particles wave equation like photons or pions or alpha nuclei that do not obey Pauli exclusion.
\end{itemize}

\noindent\textbf{Root Cause:} The fundamental obstacle is the inherent \textbf{asymmetric treatment of space and time} in quantum mechanics. \textbf{Space is an operator} ($\hat{\mathbf{X}}$), while \textbf{time is a parameter} ($t$). This asymmetry conflicts with Special Relativity's requirement of spacetime symmetry.

\noindent\textbf{Resolution Paths:} To reconcile QM with relativity, space and time must be put on an equal footing:
\begin{enumerate}
    \item \textbf{Promote Time to an Operator}: Use proper time ($\tau$) as the evolution parameter (complex, leads to string theory).
    \item \textbf{Demote Position to a Label}: \textbf{This is Quantum Field Theory}. Associate operators with spacetime points. Both $\mathbf{x}$ and $t$ become labels on operators, establishing equal footing.
\end{enumerate}

\noindent\textbf{QFT Framework:} Non-relativistic QM for multiple particles can be rewritten as a field theory:
\begin{itemize}
    \item Introduce field operators: annihilation $a(\mathbf{x})$ and creation $a^\dagger(\mathbf{x})$
    \item (Anti-)commutation relations define statistics:
    \begin{align*}
        \textbf{Bosons:} \quad & [a(\mathbf{x}), a(\mathbf{x}')] = 0, \quad [a^\dagger(\mathbf{x}), a^\dagger(\mathbf{x}')] = 0 \\
        \textbf{Fermions:} \quad & \{a(\mathbf{x}), a(\mathbf{x}')\} = 0, \quad \{a^\dagger(\mathbf{x}), a^\dagger(\mathbf{x}')\} = 0
    \end{align*}
    \item Hamiltonian in terms of fields:
    \begin{align*}
        H = \int d^3x\, a^\dagger(\mathbf{x})\left(-\frac{\hbar^2}{2m}\nabla^2 + U(\mathbf{x})\right)a(\mathbf{x}) \\
        + \frac{1}{2}\int d^3x\, d^3y\, V(\mathbf{x}-\mathbf{y})a^\dagger(\mathbf{x})a^\dagger(\mathbf{y})a(\mathbf{y})a(\mathbf{x})
    \end{align*}
    \item $n$-particle state from vacuum $\ket{0}$ ($a(\mathbf{x})\ket{0} = 0$):
    \begin{equation}
    \ket{\psi,t} = \int d^3x_1 \ldots d^3x_n\, \psi(\mathbf{x}_1, \ldots, \mathbf{x}_n; t)\, a^\dagger(\mathbf{x}_1) \ldots a^\dagger(\mathbf{x}_n) \ket{0}
    \end{equation}
    \item \textbf{Key advantage:} Variable particle number through terms like
    \begin{equation}
    \Delta H \propto \int d^3x\, \left[ a^\dagger(\mathbf{x})a^2(\mathbf{x}) + \text{h.c.} \right]
    \end{equation}
    which describes creation/annihilation processes.
\end{itemize}

\noindent\textbf{Summary:} Quantum Field Theory is the inevitable framework consistent with both QM and SR, where particles are excitations of fields and particle number is dynamical.

This section argues that \textbf{Quantum Field Theory is not merely an advanced topic but the inevitable framework} required to construct a theory that is consistent with both the principles of quantum mechanics and the laws of special relativity.
\subsection{Lorentz Invariance}

\section{Spin One Half}

\end{document}